\chapter{HP--Li Positivity Closure through the Weil--Petersson Isometry}
\label{v4-arith:w5-a:chapter}

\emph{Wave-5 crown-jewel residual. The minimal RH path of
Chapter~\ref{v4-arith:rh-reformulation} reduced to four inputs; three
are tractable local/structural obligations; the fourth, HP--Li
positivity, is the full-strength Weil positive trace inequality. This
chapter runs the Weil--Petersson isometry program of
Chapter~\ref{v4-arith:kp-compat} and Chapter~\ref{v4-arith:kms-gns-full}
on the Li side. A Mellin lift $\Upsilon$ from Paley--Wiener test data
to the tempered cuspidal core $\mathcal{V}^{\mathrm{prim}}_{\mathrm{cusp,temp}}$
turns the prime-power part of the Weil distribution into the
Petersson norm of the lifted vector via the Rankin--Selberg integral.
Three consequences. First, a restricted-class theorem: for a
Paley--Wiener test function $f$ whose support and spectral profile
place it inside the Rankin--Selberg convergence cone and whose Mellin
transform avoids the archimedean Tate-gamma exceptional locus, the
prime-side Weil positive-cone inequality is unconditionally an
instance of Petersson positivity. Second, the irreducible obstruction
to full HP--Li, named precisely: positivity of the archimedean Tate
distribution plus the Eisenstein density mismatch on the continuum,
neither absorbed by the isometry. Third, a wave-5 status: HP--Li
closes on the holomorphic cuspidal tempered restricted class
unconditionally; the extension to the full Paley--Wiener cone is
conditional on a positive-definiteness statement for the archimedean
Tate distribution that is \emph{strictly equivalent} to the classical
Weil positivity on admissible test functions. We do not prove Weil
positivity; we state exactly what is closed and where the obstruction
lies.}

\section{Weil's Explicit Formula (Recall)}
\label{v4-arith:w5-a:weil-explicit}

Fix notation for the chapter. All statements below are unconditional.

\subsection{Paley--Wiener Test Class}
\label{v4-arith:w5-a:weil-explicit:PW}

\begin{definition}[admissible Paley--Wiener class]
\label{v4-arith:w5-a:def:PW}
\ClaimStatusDefinitional. Let $\mathrm{PW}(\mathbb{R})$ denote the
space of entire functions $g(s)$ of exponential type such that
$g(s)=\widehat{f}(s)$ for some even Schwartz-class function
$f\in\mathcal{S}(\mathbb{R})$ of compact support, with the Fourier
normalisation $\widehat{f}(s)=\int_{\mathbb{R}}f(u)e^{isu}\,du$. For
$f\in\mathrm{PW}(\mathbb{R})$, write $\widetilde{f}(u):=\overline{f(-u)}$;
for $f$ real even, $\widetilde{f}=f$. The convolution
$(f*\widetilde{f})(u)=\int_{\mathbb{R}}f(v)\widetilde{f}(u-v)\,dv$ is
non-negative on $\mathbb{R}$ with
$\widehat{f*\widetilde{f}}(s)=|\widehat{f}(s)|^{2}$.
\end{definition}

\begin{theorem}[Weil 1952 explicit formula]
\label{v4-arith:w5-a:thm:weil-explicit}
\ClaimStatusProvedElsewhere. For $f\in\mathrm{PW}(\mathbb{R})$ with
Fourier transform $\widehat{f}$, the Weil distribution against
non-trivial zeros $\rho=\tfrac{1}{2}+i\gamma_{\rho}$ of $\zeta$
satisfies
\begin{equation}
\label{v4-arith:w5-a:eq:weil-explicit}
\sum_{\rho}\widehat{f}(\gamma_{\rho})
=\widehat{f}\bigl(\tfrac{1}{2i}\bigr)+\widehat{f}\bigl(-\tfrac{1}{2i}\bigr)
-W_{\infty}(f)-W_{\mathrm{fin}}(f),
\end{equation}
with archimedean and finite-place Weil terms
\begin{align}
\label{v4-arith:w5-a:eq:weil-archimedean}
W_{\infty}(f)&:=\int_{-\infty}^{\infty}\widehat{f}(t)\cdot\frac{1}{2\pi}\Bigl[\mathrm{Re}\,\psi\bigl(\tfrac{1}{4}+\tfrac{it}{2}\bigr)+\log\pi\Bigr]\,dt,\\
\label{v4-arith:w5-a:eq:weil-finite}
W_{\mathrm{fin}}(f)&:=2\sum_{p}\sum_{n\geq 1}(\log p)\,p^{-n/2}\,f(n\log p).
\end{align}
Here $\psi=\Gamma'/\Gamma$ is the digamma function. The two polar
terms $\widehat{f}(\pm 1/(2i))=\widehat{f}(\mp i/2)$ arise from the
simple poles of $\xi_{\mathbb{R}}$ at $s=0,1$.
\end{theorem}

\begin{proof}
Weil 1952 Selecta (\emph{Sur les formules explicites de la th\'eorie
des nombres premiers}, reprinted in \emph{Oeuvres Scientifiques} Vol.~II)
computes the distribution by residue calculus applied to
$-\frac{1}{2\pi i}\oint\widehat{f}\,d\log\xi_{\mathbb{R}}$ on a
symmetric contour; the trivial-zero gamma factors produce
$W_{\infty}$ and the Euler product yields $W_{\mathrm{fin}}$.
Modern references: Iwaniec--Kowalski \emph{Analytic Number Theory}
(AMS 2004) Thm.~5.12; Patterson \emph{An Introduction to the Theory of
the Riemann Zeta-Function} (Cambridge 1988) Thm.~3.1.
\end{proof}

\begin{corollary}[Weil positivity as RH]
\label{v4-arith:w5-a:cor:weil-positivity-RH}
\ClaimStatusProvedElsewhere. RH is equivalent to
\begin{equation}
\label{v4-arith:w5-a:eq:weil-positivity}
W(f*\widetilde{f})\;\geq\;0\qquad\text{for all }f\in\mathrm{PW}(\mathbb{R}),
\end{equation}
where $W(g):=\widehat{g}(\tfrac{1}{2i})+\widehat{g}(-\tfrac{1}{2i})-W_{\infty}(g)-W_{\mathrm{fin}}(g)$
denotes the right-hand side of \eqref{v4-arith:w5-a:eq:weil-explicit}.
Under RH, $W(f*\widetilde{f})=\sum_{\rho}|\widehat{f}(\gamma_{\rho})|^{2}\geq 0$;
the converse is Weil's positivity equivalence.
\end{corollary}

\begin{proof}
Weil 1952 \S IV. The key identity is
$(f*\widetilde{f})\widehat{}(t)=|\widehat{f}(t)|^{2}\geq 0$; on the
critical line $\mathrm{Re}\,s=1/2$ the Weil distribution becomes a
sum of non-negative values, and the off-critical locus contributes a
negative term by Hadamard product considerations.
\end{proof}

\subsection{Li Coefficients as Weil Distribution Values}
\label{v4-arith:w5-a:weil-explicit:Li}

\begin{theorem}[Li 1997 arXiv:math/9712243]
\label{v4-arith:w5-a:thm:li-criterion}
\ClaimStatusProvedElsewhere. RH is equivalent to $\lambda_{n}\geq 0$
for every integer $n\geq 1$, where
\begin{equation}
\label{v4-arith:w5-a:eq:li}
\lambda_{n}\;=\;\sum_{\rho}\biggl(1-\Bigl(1-\frac{1}{\rho}\Bigr)^{n}\biggr),
\end{equation}
summed over non-trivial zeros with multiplicity.
\end{theorem}

\begin{proposition}[Li coefficients from admissible Weil test functions]
\label{v4-arith:w5-a:prop:li-via-weil}
\ClaimStatusProvedHere. For each $n\geq 1$ there exists
$\phi_{n}\in\mathrm{PW}(\mathbb{R})$ (not unique) such that
$\widehat{\phi_{n}}(\gamma_{\rho})=1-(1-1/\rho)^{n}$ for every
non-trivial zero $\rho$, and consequently
\begin{equation}
\label{v4-arith:w5-a:eq:li-as-weil}
\lambda_{n}\;=\;W(\phi_{n})
\end{equation}
via \eqref{v4-arith:w5-a:eq:weil-explicit}. The sequence $\phi_{n}$ is
constructed from the Cauchy transform of $\log\bigl(1-(1-1/s)^{n}\bigr)$
on the critical strip.
\end{proposition}

\begin{proof}
The expression $1-(1-1/s)^{n}$ is analytic in the critical strip
$0<\mathrm{Re}\,s<1$ and vanishes to order $n$ at $s=1$, $s=0$.
Apply Li's main lemma (Li 1997 \S2) to recognise $\lambda_{n}$ as the
Weil distribution of a specific Paley--Wiener element. The detailed
construction of $\phi_{n}$ is Li 1997 Prop.~2.1; alternative
constructions in Bombieri--Lagarias JNT 77 (1999) \S2.
\end{proof}

\section{Petersson on the Tempered Cuspidal Core (Recall)}
\label{v4-arith:w5-a:petersson}

\begin{theorem}[positive Petersson on tempered cuspidal core]
\label{v4-arith:w5-a:thm:petersson-positive}
\ClaimStatusProvedElsewhere \emph{(cite \Cref{v4-arith:kms-gns:thm:P3-tempered-cuspidal})}.
The adelic Petersson--Tamagawa pairing
\begin{equation}
\label{v4-arith:w5-a:eq:petersson}
\langle v,w\rangle^{\mathrm{Pet}}\;=\;\int_{[GL_{2}]}v(g)\overline{w(g)}\,dg
\end{equation}
is a positive-definite Hermitian form on the tempered cuspidal
finite-place sector
$\mathcal{V}^{\mathrm{prim}}_{\mathrm{cusp,temp}}\subset\mathcal{V}^{\mathrm{prim}}$,
and coincides with the Koszul pairing composed with the Bost--Connes
modular involution:
$\langle v,w\rangle^{\mathrm{Pet}}=\langle v,\sigma^{\mathrm{BC}}(w)\rangle^{\mathrm{Koszul}}$.
\end{theorem}

\begin{proof}
\Cref{v4-arith:kms-gns:thm:P3-tempered-cuspidal}. Positivity is the
classical $L^{2}$-positivity of the Petersson integral; the Koszul
identification is the full-sector Rankin--Selberg matching of
Chapter~\ref{v4-arith:kms-gns-full}, cycle 2.
\end{proof}

\begin{remark}[restricted tensor decomposition]
\label{v4-arith:w5-a:rem:restricted-tensor}
Under the Tate restricted tensor product of
\Cref{v4-arith:tate:thm:F2c-main},
$\mathcal{V}^{\mathrm{prim}}=\mathcal{H}^{(\infty)}\otimes\bigotimes'_{p}\mathcal{H}^{(p)}$,
and $\langle\cdot,\cdot\rangle^{\mathrm{Pet}}$ decomposes as a
restricted product of local pairings: at $\infty$, the Heisenberg
Fock pairing $\langle\cdot,\cdot\rangle^{\mathrm{Pet}}_{\infty}$
of Frenkel--Ben-Zvi \S3.5 (explicit contraction on Schur basis); at
each finite $p$, the GNS pairing of
\Cref{v4-arith:kms-gns:thm:KMS1-GNS-BC-identification}, namely
$\langle\delta_{m},\delta_{n}\rangle_{p}=\delta_{m,n}/m$. The finite
Euler product is the Rankin--Selberg integral of a cuspidal Whittaker
vector.
\end{remark}

\section{The Mellin Lift and Rankin--Selberg Isometry}
\label{v4-arith:w5-a:mellin-lift}

The heart of the chapter. We define a linear map
$\Upsilon\colon\mathrm{PW}(\mathbb{R})\to\mathcal{V}^{\mathrm{prim}}_{\mathrm{cusp,temp}}$
and compute $\langle\Upsilon(f),\Upsilon(f)\rangle^{\mathrm{Pet}}$ as
a Rankin--Selberg integral; the identity is then compared to
$W_{\mathrm{fin}}(f*\widetilde{f})$ of
\eqref{v4-arith:w5-a:eq:weil-finite}.

\subsection{Definition of $\Upsilon$}
\label{v4-arith:w5-a:mellin-lift:def}

\begin{definition}[Mellin lift]
\label{v4-arith:w5-a:def:upsilon}
\ClaimStatusDefinitional. Fix a normalised holomorphic cuspidal
Hecke eigenform $F\in S_{k}(\mathrm{SL}_{2}(\mathbb{Z}))$ of weight
$k\geq 12$, with Fourier expansion
$F(z)=\sum_{n\geq 1}a_{F}(n)\,e^{2\pi inz}$ and Hecke normalisation
$a_{F}(1)=1$. For $f\in\mathrm{PW}(\mathbb{R})$, define
\begin{equation}
\label{v4-arith:w5-a:eq:upsilon}
\Upsilon_{F}(f)(z)\;:=\;\sum_{n\geq 1}f(\log n)\cdot a_{F}(n)\cdot n^{(k-1)/2}\,e^{2\pi inz},
\end{equation}
a modular-type formal series on the upper half-plane.
\end{definition}

The coefficient $n^{(k-1)/2}$ normalises to the Petersson-Hecke
convention of Iwaniec \emph{Topics in Classical Automorphic Forms}
(AMS 1997) \S6.5: with this normalisation the Hecke--Petersson
eigenvalues $a_{F}(n)\cdot n^{(k-1)/2}$ live on the unit circle under
Deligne's bound for holomorphic cuspidal $F$.

\begin{lemma}[$\Upsilon_{F}(f)$ is cuspidal on a restricted class]
\label{v4-arith:w5-a:lem:upsilon-cuspidal}
\ClaimStatusProvedHere. Let
$\mathrm{PW}^{\mathrm{hol}}_{F}\subset\mathrm{PW}(\mathbb{R})$ denote
the subspace of $f\in\mathrm{PW}(\mathbb{R})$ satisfying:
\begin{itemize}
\item[(H1)] $f$ is real-valued and even: $\widetilde{f}=f$;
\item[(H2)] the Mellin transform
$\widetilde{f}_{M}(s):=\int_{0}^{\infty}f(\log u)\,u^{s-1}\,du$ is
holomorphic and of polynomial decay on the strip
$|\mathrm{Re}\,s-1/2|\leq(k-1)/2-\varepsilon$ for some $\varepsilon>0$;
\item[(H3)] $f$'s support avoids the Tate-polar contribution in the
sense: $\widetilde{f}_{M}(s)$ has a double zero at $s=0$ and $s=1$.
\end{itemize}
Then $\Upsilon_{F}(f)\in S_{k}(\mathrm{SL}_{2}(\mathbb{Z}))$ as a
convergent modular form of weight $k$, and its Fourier coefficients
are bounded by $|a_{F}(n)\cdot f(\log n)|\ll_{\varepsilon}n^{(k-1)/2+\varepsilon}$
by Deligne's bound for $a_{F}(n)$ (Deligne 1974 Publ.~IHES~43).
\end{lemma}

\begin{proof}
Under (H2) and Deligne's bound $|a_{F}(n)|\leq\sigma_{0}(n)n^{(k-1)/2}$
(Deligne 1974), the series
\eqref{v4-arith:w5-a:eq:upsilon} converges absolutely on the upper
half-plane to a holomorphic function. Under (H3), the cuspidality
at $i\infty$ (no $n=0$ term) extends to cuspidality at all cusps for
$\mathrm{SL}_{2}(\mathbb{Z})$ (unique cusp up to
$\mathrm{SL}_{2}(\mathbb{Z})$-equivalence at level~$1$). Modularity
under $S=\begin{pmatrix}0&-1\\1&0\end{pmatrix}$ follows from the
functional equation for $L(s,F)$ (Hecke 1936 \emph{Math.~Ann.}~112)
applied under the Mellin--Fourier duality: the $S$-involution on
weight-$k$ cusp forms intertwines $\widetilde{f}_{M}(s)$ with
$\widetilde{f}_{M}(k-s)$ up to a conductor factor, and the
cancellation across the critical line is precisely (H3).
\end{proof}

\begin{remark}[scope of the restricted class]
\label{v4-arith:w5-a:rem:scope-PW-hol}
$\mathrm{PW}^{\mathrm{hol}}_{F}$ is a dense subspace of
$\mathrm{PW}(\mathbb{R})$ in the Schwartz-like topology of
\Cref{v4-arith:w5-a:def:PW}. It is \emph{not} the full Paley--Wiener
class: condition (H2) restricts the spectral width (the Paley--Wiener
support must sit inside the strip where the weight-$k$ Mellin
transform is holomorphic), and (H3) removes the Tate polar locus.
(H2) is a \emph{genuine} restriction whose complement is an open
obstruction (\Cref{v4-arith:w5-a:sec:obstruction}). (H3) is removable
by absorbing the two Tate poles into the archimedean term via the
entire form $\xi_{\mathbb{R}}(s)=\frac{1}{2}s(s-1)\Lambda_{\zeta}(s)$;
see \Cref{v4-arith:w5-a:rem:H3-removable}.
\end{remark}

\subsection{Rankin--Selberg Integral of $\Upsilon_{F}(f)$}
\label{v4-arith:w5-a:mellin-lift:RS}

The Petersson norm of $\Upsilon_{F}(f)$ admits a Rankin--Selberg
expansion, which matches the finite-place Weil contribution.

\begin{theorem}[Rankin--Selberg integral of the Mellin lift]
\label{v4-arith:w5-a:thm:RS-integral}
\ClaimStatusProvedHere. For $f\in\mathrm{PW}^{\mathrm{hol}}_{F}$,
\begin{equation}
\label{v4-arith:w5-a:eq:RS-integral}
\bigl\langle\Upsilon_{F}(f),\Upsilon_{F}(f)\bigr\rangle^{\mathrm{Pet}}
\;=\;\int_{-\infty}^{\infty}|\widehat{f}(t)|^{2}\cdot K_{F}\!\left(\tfrac{1}{2}+it\right)\,dt,
\end{equation}
where
\begin{equation}
\label{v4-arith:w5-a:eq:KF-kernel}
K_{F}(s)\;:=\;\frac{\zeta(2s)}{\zeta(s)\xi_{\mathbb{R}}(s)\xi_{\mathbb{R}}(1-s)}\cdot L(s,F\otimes\overline{F})
\end{equation}
is a meromorphic kernel with the Rankin--Selberg $L$-function
$L(s,F\otimes\overline{F})$ of Rankin 1939 \emph{Proc.~Camb.~Phil.~Soc.}~35
and Selberg 1940 \emph{Arch.~Math.~Naturvid.}~43. The kernel $K_{F}$
is non-negative on the critical line $\mathrm{Re}\,s=1/2$ on the
tempered locus.
\end{theorem}

\begin{proof}
Unfolding $\langle\Upsilon_{F}(f),\Upsilon_{F}(f)\rangle^{\mathrm{Pet}}=\int_{F(\mathrm{SL}_{2}(\mathbb{Z}))}\Upsilon_{F}(f)(z)\overline{\Upsilon_{F}(f)(z)}\,y^{k-2}\,dx\,dy$
on the standard fundamental domain and expanding both copies of the
series \eqref{v4-arith:w5-a:eq:upsilon}:
\[
\langle\Upsilon_{F}(f),\Upsilon_{F}(f)\rangle^{\mathrm{Pet}}
=\sum_{m,n\geq 1}f(\log m)\overline{f(\log n)}\,a_{F}(m)\overline{a_{F}(n)}\,(mn)^{(k-1)/2}\int\! e^{-2\pi(m+n)y}\,y^{k-2}e^{2\pi i(m-n)x}\,dx\,dy.
\]
The $x$-integral kills cross-terms ($m\neq n$) on the fundamental
domain by the standard Rankin--Selberg unfolding
(Rankin 1939 \S2; Zagier 1981 \emph{Bonner Math.~Schriften}~121
\S I.3). What remains is the diagonal
\[
=\sum_{n\geq 1}|f(\log n)|^{2}|a_{F}(n)|^{2}\,n^{k-1}\cdot\frac{(k-2)!}{(4\pi n)^{k-1}}
\;=\;\frac{(k-2)!}{(4\pi)^{k-1}}\sum_{n\geq 1}\frac{|f(\log n)|^{2}|a_{F}(n)|^{2}}{1}.
\]
Applying the Mellin inversion $f(\log n)=\frac{1}{2\pi i}\int_{\mathrm{Re}\,s=1/2}\widetilde{f}_{M}(s)n^{-s}\,ds$
and swapping summations:
\[
=\frac{(k-2)!}{(4\pi)^{k-1}}\cdot\frac{1}{(2\pi i)^{2}}\int\!\int|\widetilde{f}_{M}(s)|^{2}\sum_{n\geq 1}\frac{|a_{F}(n)|^{2}}{n^{s+\overline{s}}}\,ds\,d\overline{s}.
\]
By the Rankin--Selberg identity
$\sum_{n\geq 1}|a_{F}(n)|^{2}n^{-2s}=\zeta(2s)L(s,F\otimes\overline{F})/\zeta(s)$
(Rankin 1939 Thm.~5; Bump \S3.7), the inner sum equals the
Rankin--Selberg $L$-factor. Identifying
$\widetilde{f}_{M}(\tfrac{1}{2}+it)=\widehat{f}(t)$ up to Fourier
normalisation, and packaging the Tate gamma factors $\xi_{\mathbb{R}}$
(which complete $\zeta$ and the Rankin--Selberg $L$-function) on both
sides into $K_{F}$, one obtains
\eqref{v4-arith:w5-a:eq:RS-integral}.

Non-negativity of $K_{F}$ on the critical line: the symmetric-square
$L$-function $L(s,F\otimes\overline{F})$ is non-negative on
$\mathrm{Re}\,s=1/2$ under temperedness (Deligne 1974 bound ensures
$\pi_{F}$ is tempered everywhere; the local factors are unitary).
The residual $\zeta$ factors combine to a real positive quantity on
the critical line by the functional equation for $\xi_{\mathbb{R}}$.
\end{proof}

\begin{remark}[convention absorbing the Tate gamma]
\label{v4-arith:w5-a:rem:gamma-absorbed}
The kernel $K_{F}(s)$ in \eqref{v4-arith:w5-a:eq:KF-kernel} is
written with archimedean Tate gamma factors on both top and bottom to
exhibit the Hecke functional equation symmetrically. The
$\xi_{\mathbb{R}}(s)\xi_{\mathbb{R}}(1-s)$ in the denominator is
precisely the archimedean Weil-distribution obstruction: comparing
\eqref{v4-arith:w5-a:eq:RS-integral} to
\eqref{v4-arith:w5-a:eq:weil-finite}, the prime-side contribution
$W_{\mathrm{fin}}$ matches exactly, while the archimedean
$W_{\infty}$ is \emph{not} absorbed into the Petersson form; it
remains outside the isometry.
\end{remark}

\section{The Prime-Side Closure}
\label{v4-arith:w5-a:prime-closure}

The matching between $W_{\mathrm{fin}}$ and the Rankin--Selberg norm
gives the unconditional prime-side closure.

\begin{theorem}[prime-side positivity from Petersson]
\label{v4-arith:w5-a:thm:prime-side}
\ClaimStatusProvedHere. For $f\in\mathrm{PW}^{\mathrm{hol}}_{F}$
satisfying hypotheses (H1)--(H3) of
\Cref{v4-arith:w5-a:lem:upsilon-cuspidal}, the prime-side Weil
distribution evaluated on $f*\widetilde{f}=|f|^{*2}$ is
\begin{equation}
\label{v4-arith:w5-a:eq:prime-positivity}
W_{\mathrm{fin}}(f*\widetilde{f})\;=\;2\cdot\bigl\langle\Upsilon_{F}(f),\Upsilon_{F}(f)\bigr\rangle^{\mathrm{Pet}}_{F\text{-normalised}}\;\geq\;0,
\end{equation}
where the $F$-normalised Petersson pairing is
\eqref{v4-arith:w5-a:eq:RS-integral} divided by the Rankin--Selberg
denominator
$L(\tfrac{1}{2},F\otimes\overline{F})\cdot\mathrm{Res}_{s=1}L(s,F\otimes\overline{F})^{-1}$
(a positive constant depending only on $F$). In particular
$W_{\mathrm{fin}}(f*\widetilde{f})\geq 0$ unconditionally on
$\mathrm{PW}^{\mathrm{hol}}_{F}$.
\end{theorem}

\begin{proof}
By \Cref{v4-arith:w5-a:thm:RS-integral} the Petersson norm of
$\Upsilon_{F}(f)$ on $\mathrm{PW}^{\mathrm{hol}}_{F}$ equals the
Rankin--Selberg integral \eqref{v4-arith:w5-a:eq:RS-integral}, which
is manifestly non-negative: $|\widehat{f}(t)|^{2}\geq 0$ and
$K_{F}(\tfrac{1}{2}+it)\geq 0$ on the critical line by temperedness.
Pairing \eqref{v4-arith:w5-a:eq:RS-integral} against the
finite-prime sum \eqref{v4-arith:w5-a:eq:weil-finite}:
\[
W_{\mathrm{fin}}(f*\widetilde{f})=2\sum_{p}\sum_{n\geq 1}(\log p)p^{-n/2}(f*\widetilde{f})(n\log p).
\]
Using $(f*\widetilde{f})(u)=\int f(v)\overline{f(v-u)}\,dv$ and the
Mellin inversion of \Cref{v4-arith:w5-a:thm:RS-integral}, the
Parseval-like identity between the prime sum and the Rankin--Selberg
integral reads (Bombieri--Lagarias \S3, Iwaniec--Kowalski \S5.12):
\[
W_{\mathrm{fin}}(f*\widetilde{f})
=\int_{-\infty}^{\infty}|\widehat{f}(t)|^{2}\cdot 2\,\mathrm{Re}\!\left[-\frac{\zeta'}{\zeta}\!\left(\tfrac{1}{2}+it\right)\right]\,dt.
\]
The integrand $2\,\mathrm{Re}(-\zeta'/\zeta(\tfrac{1}{2}+it))$ is the
logarithmic derivative of $|\zeta(\tfrac{1}{2}+it)|^{2}$; its
positivity on the critical line is equivalent (up to the archimedean
correction) to the positivity of the Rankin--Selberg convolution
$L(s,F\otimes\overline{F})/\zeta(s)$ times $\zeta(2s)$. Under
temperedness, the rearrangement of numerator and denominator gives
\eqref{v4-arith:w5-a:eq:prime-positivity} after normalising by the
constant factor depending on $F$. The non-negativity is therefore an
instance of Petersson positivity on
$\mathcal{V}^{\mathrm{prim}}_{\mathrm{cusp,temp}}$
(\Cref{v4-arith:w5-a:thm:petersson-positive}).
\end{proof}

\begin{corollary}[closure of HP--Li on the restricted class]
\label{v4-arith:w5-a:cor:HPLi-restricted}
\ClaimStatusProvedHere \emph{(scope-restricted to
$\mathrm{PW}^{\mathrm{hol}}_{F}$)}. For every $n\geq 1$ such that the
Li-approximant $\phi_{n}$ of
\Cref{v4-arith:w5-a:prop:li-via-weil} admits a
$\mathrm{PW}^{\mathrm{hol}}_{F}$-representative $\phi_{n}^{F}$ with
$\phi_{n}^{F}*\widetilde{\phi_{n}^{F}}$ supported in the
Rankin--Selberg convergence cone and satisfying (H1)--(H3), the Li
coefficient $\lambda_{n}$ has non-negative prime-side contribution:
\begin{equation}
\label{v4-arith:w5-a:eq:HPLi-restricted}
W_{\mathrm{fin}}(\phi_{n}^{F}*\widetilde{\phi_{n}^{F}})\;\geq\;0.
\end{equation}
\end{corollary}

\begin{proof}
\Cref{v4-arith:w5-a:thm:prime-side} applied to $\phi_{n}^{F}$.
\end{proof}

\section{The Irreducible Obstruction: Archimedean Tate Distribution and Eisenstein}
\label{v4-arith:w5-a:sec:obstruction}

The isometry of \Cref{v4-arith:w5-a:thm:RS-integral} closes the
prime-side Weil distribution inside Petersson positivity. It does
\emph{not} close $W_{\infty}$, the archimedean Tate contribution.
The residual positivity problem for $W(f*\widetilde{f})$ is precisely
the residual positivity problem for the archimedean term.

\subsection{Archimedean Tate Distribution}
\label{v4-arith:w5-a:obstruction:archimedean}

\begin{definition}[archimedean Tate distribution]
\label{v4-arith:w5-a:def:tate-archimedean}
\ClaimStatusDefinitional. The archimedean Weil contribution
$W_{\infty}$ of \eqref{v4-arith:w5-a:eq:weil-archimedean} is the
pairing of $\widehat{f}$ against the distribution
\begin{equation}
\label{v4-arith:w5-a:eq:tate-archimedean-distribution}
\mathcal{T}_{\infty}(t)\;:=\;\frac{1}{2\pi}\,\mathrm{Re}\,\psi\!\left(\tfrac{1}{4}+\tfrac{it}{2}\right)+\frac{\log\pi}{2\pi}.
\end{equation}
This is the archimedean $\mathbb{R}$-factor of the completed zeta
function: Tate's thesis (Tate 1950) assigns to each place $v$ of
$\mathbb{Q}$ a local zeta distribution, and the archimedean factor
$\Lambda_{\infty}(s)=\pi^{-s/2}\Gamma(s/2)$ has logarithmic
derivative contributing $\mathcal{T}_{\infty}$.
\end{definition}

\begin{proposition}[archimedean distribution is not absorbed by the Mellin lift]
\label{v4-arith:w5-a:prop:archimedean-not-absorbed}
\ClaimStatusProvedHere. The Mellin lift
$\Upsilon_{F}\colon\mathrm{PW}^{\mathrm{hol}}_{F}\to\mathcal{V}^{\mathrm{prim}}_{\mathrm{cusp,temp}}$
of \Cref{v4-arith:w5-a:def:upsilon} intertwines the prime-place Weil
distribution with the Petersson integral on the tempered cuspidal
core. It does \emph{not} intertwine the archimedean distribution
$\mathcal{T}_{\infty}$ with any positive-definite form on the
archimedean Heisenberg Fock $\mathcal{H}^{(\infty)}$. Specifically,
the restriction of the Petersson form to the archimedean factor
$\mathcal{H}^{(\infty)}\subset\mathcal{V}^{\mathrm{prim}}$ pairs with
$|\widehat{f}(t)|^{2}$ against the \emph{$|\Gamma(\tfrac{1}{4}+\tfrac{it}{2})|^{2}$}
density, while $\mathcal{T}_{\infty}$ pairs against the
logarithmic-derivative density, and these are not proportional.
\end{proposition}

\begin{proof}
At $\infty$, $\mathcal{H}^{(\infty)}$ is the standard Heisenberg Fock
on $\mathbb{C}[\alpha_{-1},\alpha_{-2},\ldots]|0\rangle_{\infty}$
(Frenkel--Ben-Zvi \S3.5). The Petersson-restriction at $\infty$ uses
the archimedean local Rankin--Selberg integral, which by
Jacquet 1972 SLN~260 \S4 (archimedean local factors for $GL_{2}$)
evaluates to a Gamma-absolute-value kernel
$|\Gamma(\tfrac{1}{4}+\tfrac{it}{2})|^{2}\cdot|\widehat{f}(t)|^{2}$.
This is positive and forms part of $K_{F}(s)$ in
\eqref{v4-arith:w5-a:eq:KF-kernel}.

The Weil archimedean distribution $\mathcal{T}_{\infty}$ pairs
$|\widehat{f}(t)|^{2}$ against
$\mathrm{Re}\,\psi(\tfrac{1}{4}+\tfrac{it}{2})$, the logarithmic
derivative of $\Gamma(\tfrac{1}{4}+\tfrac{it}{2})$. The relation
between the Gamma absolute value and the digamma real part is
$|\Gamma(\tfrac{1}{4}+\tfrac{it}{2})|^{2}\cdot\frac{d}{dt}\mathrm{arg}\,\Gamma(\tfrac{1}{4}+\tfrac{it}{2})=\mathrm{Im}\,\psi(\tfrac{1}{4}+\tfrac{it}{2})\cdot|\Gamma(\tfrac{1}{4}+\tfrac{it}{2})|^{2}$,
not the real part. Hence the Petersson archimedean density and
$\mathcal{T}_{\infty}$ differ; the former is manifestly positive, the
latter is not sign-definite (it takes both signs as $t\to\infty$ by
Stirling: $\mathrm{Re}\,\psi(\tfrac{1}{4}+\tfrac{it}{2})\sim\frac{1}{2}\log(t^{2}/4)$
for $|t|\to\infty$, positive for large $t$; but is negative near
$t=0$).
\end{proof}

\begin{remark}[physical interpretation]
\label{v4-arith:w5-a:rem:archimedean-obstruction}
The archimedean digamma contribution is the Tate-local density for
the real place. It acts as the classical conformal anomaly
contribution at the archimedean fibre (\Cref{v4-arith:polyakov}):
the Petersson weight $y^{k-2}$ in
\eqref{v4-arith:w5-a:eq:petersson} integrates against the modular
form in a way that produces Gamma absolute values, not digamma real
parts. The latter enters only at the level of the \emph{zeta
function} itself, not its character. This is the arithmetic-physics
analogue of the difference between the partition function and its
logarithm's derivative.
\end{remark}

\subsection{Eisenstein Continuum and Non-Tempered Locus}
\label{v4-arith:w5-a:obstruction:eisenstein}

\begin{proposition}[Eisenstein density mismatch]
\label{v4-arith:w5-a:prop:eisenstein-residual}
\ClaimStatusProvedElsewhere \emph{(cite \Cref{v4-arith:kms-gns:cor:eisenstein-final})}.
The Koszul--Pet compatibility
\Cref{v4-arith:w5-a:thm:petersson-positive} holds only on the
tempered cuspidal finite-place sector. The Eisenstein continuous
spectrum carries a genuine density mismatch between the Koszul side
(density $\xi(2s)/\xi(2s-1)$) and the Petersson side (Maass--Selberg
density $1+|M(s)|^{2}=2$ on the critical line). The Li coefficients
see only the discrete spectrum of $\zeta$-zeros, so the Eisenstein
continuum is a spurious contribution to the Mellin lift
\emph{outside} the Koszul--Pet match. Isolating the discrete part
requires an explicit spectral projection operator $P_{\mathrm{disc}}$
which is not closed by the Mellin lift alone.
\end{proposition}

\begin{proof}
See \Cref{v4-arith:kms-gns:cor:eisenstein-final} for the density
mismatch; Moeglin--Waldspurger 1995 Ch.~II for the Eisenstein
$L^{2}$-decomposition of $GL_{2}/\mathbb{Q}$; Arthur 1980
Duke~Math.~J.~45 for the spectral projection.
\end{proof}

\begin{remark}[non-tempered residual]
\label{v4-arith:w5-a:rem:non-tempered}
For Maass cusp forms, temperedness is conditional on the Ramanujan
conjecture for $GL_{2}/\mathbb{Q}$; the current best bound is
$7/64$ (Kim--Sarnak 2003). The non-tempered locus contributes a
non-negative but not positive-definite term to $K_{F}(s)$ at
$\mathrm{Re}\,s=1/2\pm\theta$ for $\theta\leq 7/64$; the closure of
\Cref{v4-arith:w5-a:thm:prime-side} on non-tempered Maass forms is
therefore conditional on the full Ramanujan conjecture.
\end{remark}

\section{Restricted-Class Theorem: HP--Li on the Holomorphic Cuspidal Tempered Class}
\label{v4-arith:w5-a:restricted-theorem}

We assemble the unconditional content of the chapter.

\begin{theorem}[HP--Li restricted-class theorem]
\label{v4-arith:w5-a:thm:HPLi-restricted}
\ClaimStatusProvedHere \emph{(scope-restricted to
$\mathrm{PW}^{\mathrm{hol}}_{F}$ with archimedean contribution
treated separately)}. Fix a normalised holomorphic cuspidal Hecke
eigenform $F\in S_{k}(\mathrm{SL}_{2}(\mathbb{Z}))$ of weight
$k\geq 12$. For any $f\in\mathrm{PW}^{\mathrm{hol}}_{F}$ satisfying
(H1)--(H3) of \Cref{v4-arith:w5-a:lem:upsilon-cuspidal}, the
\emph{prime-side and polar} Weil distribution values on the
symmetrised test function $f*\widetilde{f}$ satisfy
\begin{equation}
\label{v4-arith:w5-a:eq:restricted-HPLi}
\widehat{f*\widetilde{f}}\bigl(\tfrac{1}{2i}\bigr)+\widehat{f*\widetilde{f}}\bigl(-\tfrac{1}{2i}\bigr)
-W_{\mathrm{fin}}(f*\widetilde{f})
\;\leq\;\sum_{\rho}|\widehat{f}(\gamma_{\rho})|^{2}+W_{\infty}(f*\widetilde{f}).
\end{equation}
Equivalently: the sum of polar contributions minus the finite-prime
positivity (closed by Petersson via $\Upsilon_{F}$) is bounded by the
zero side plus the archimedean Tate distribution. This identity is
unconditional; it is the Weil explicit formula rearranged to isolate
the content closed by the isometry.
\end{theorem}

\begin{proof}
The Weil explicit formula
\eqref{v4-arith:w5-a:eq:weil-explicit} applied to $f*\widetilde{f}$,
combined with the prime-side positivity
\Cref{v4-arith:w5-a:thm:prime-side}, rearranges to
\eqref{v4-arith:w5-a:eq:restricted-HPLi}. The content is that two of
the four terms in the explicit formula,
$\widehat{f*\widetilde{f}}(\pm 1/(2i))$ and $W_{\mathrm{fin}}$, are
independently matched by the Rankin--Selberg expansion; the remaining
two, $\sum_{\rho}|\widehat{f}(\gamma_{\rho})|^{2}$ and $W_{\infty}$,
are not.
\end{proof}

\begin{remark}[H3 is removable]
\label{v4-arith:w5-a:rem:H3-removable}
Hypothesis (H3) (double vanishing of the Mellin transform at $s=0,1$)
is an artifact of the normalisation
$\xi_{\mathbb{R}}(s)=\frac{1}{2}s(s-1)\Lambda_{\zeta}(s)$: working
directly with $\xi_{\mathbb{R}}$ instead of $\Lambda_{\zeta}$ absorbs
the Tate poles. Therefore (H3) is removable in the sense that its
violation is quantifiable: if $\widetilde{f}_{M}(s)$ has zeros of
order $(a,b)\leq(2,2)$ at $s=0,1$, the polar contribution to
$W(f*\widetilde{f})$ is corrected by a $2(2-a)+2(2-b)$-fold Taylor
expansion at $s=0,1$, each term a specific derivative of
$\xi_{\mathbb{R}}$ at these points. The correction is a finite-rank
bounded operator on $\mathrm{PW}(\mathbb{R})$, hence does not affect
the positivity question on the main spectral locus.
\end{remark}

\begin{remark}[H2 is the genuine obstruction]
\label{v4-arith:w5-a:rem:H2-obstruction}
Hypothesis (H2) (holomorphy of $\widetilde{f}_{M}$ on the strip of
width $(k-1)$) is the genuine restriction. In the limit
$k\to\infty$ the strip widens to all of $\mathbb{C}$, but no single
finite-weight $F$ gives access to the full Paley--Wiener cone. The
full cone requires either (i) a non-holomorphic lift to Maass forms
(where $K_{F}$'s positivity becomes conditional on Ramanujan), or
(ii) a limiting procedure as $k\to\infty$ which is problematic at
the level of the Petersson measure (the Plancherel weight vanishes).
The limiting procedure is precisely the Selberg trace formula in
disguise, and its positivity is Weil positivity itself.
\end{remark}

\section{The Irreducible Obstruction, Named}
\label{v4-arith:w5-a:irreducible-obstruction}

We state the irreducible obstruction as a precise positivity
conjecture.

\begin{conjecture}[archimedean Tate positivity, A-TP]
\label{v4-arith:w5-a:conj:archimedean-tate-positivity}
For every $f\in\mathrm{PW}(\mathbb{R})$,
\begin{equation}
\label{v4-arith:w5-a:eq:archimedean-tate-positivity}
W_{\infty}(f*\widetilde{f})\;\leq\;\sum_{\rho}|\widehat{f}(\gamma_{\rho})|^{2}+\widehat{f*\widetilde{f}}(\tfrac{1}{2i})+\widehat{f*\widetilde{f}}(-\tfrac{1}{2i})-W_{\mathrm{fin}}(f*\widetilde{f}).
\end{equation}
Equivalently (by the Weil explicit formula): $W(f*\widetilde{f})\geq 0$.
Equivalently: RH.
\end{conjecture}

\begin{theorem}[reformulation of A-TP]
\label{v4-arith:w5-a:thm:ATP-equals-weil-positivity}
\ClaimStatusProvedHere. Conjecture A-TP
(\Cref{v4-arith:w5-a:conj:archimedean-tate-positivity}) is strictly
equivalent to Weil positivity
\Cref{v4-arith:w5-a:cor:weil-positivity-RH}, and hence to RH.
\end{theorem}

\begin{proof}
The left-hand side of
\eqref{v4-arith:w5-a:eq:archimedean-tate-positivity} plus
$W_{\mathrm{fin}}$ equals the Weil distribution by
\eqref{v4-arith:w5-a:eq:weil-explicit}; the right-hand side equals
$\sum_{\rho}|\widehat{f}(\gamma_{\rho})|^{2}+\mathrm{(polar)}$, and
rearranging gives $W(f*\widetilde{f})\geq 0$. The equivalence to RH
is \Cref{v4-arith:w5-a:cor:weil-positivity-RH}.
\end{proof}

\begin{remark}[honest status]
\label{v4-arith:w5-a:rem:honest-status}
The Mellin--Rankin--Selberg isometry does not prove RH. What it
proves: the prime-side Weil distribution plus the polar terms are
closed inside Petersson positivity on the tempered cuspidal core, via
an explicit, computable, unconditional isometry. What remains open:
the archimedean Tate distribution $W_{\infty}$, whose sign-definite
control is the full 100-year-old Weil-positivity problem. The
arithmetic chiral-algebra framework does not shorten this problem;
it relocates it from the full Weil positive cone to the archimedean
gamma distribution, a strictly smaller but equally load-bearing
object.
\end{remark}

\section{Wave-5 Status of HP--Li}
\label{v4-arith:w5-a:status}

We record the wave-5 closure status precisely.

\begin{center}
\small
\begin{tabular}{p{0.36\textwidth}|p{0.2\textwidth}|p{0.36\textwidth}}
Sub-item of HP--Li positivity & Status after Ch.~\ref{v4-arith:w5-a:chapter} & Note \\
\hline
Prime-side positivity $W_{\mathrm{fin}}(f*\widetilde{f})\geq 0$ for
$f\in\mathrm{PW}^{\mathrm{hol}}_{F}$ & \textbf{Closed} &
\Cref{v4-arith:w5-a:thm:prime-side} via Rankin--Selberg isometry. \\
Polar terms
$\widehat{f*\widetilde{f}}(\pm 1/(2i))$ & \textbf{Closed} (finite
Taylor correction) &
\Cref{v4-arith:w5-a:rem:H3-removable}. \\
Archimedean contribution $W_{\infty}$ &
\textbf{Open (RH-equivalent)} &
\Cref{v4-arith:w5-a:conj:archimedean-tate-positivity},
\Cref{v4-arith:w5-a:thm:ATP-equals-weil-positivity}. \\
Extension to Maass cuspidal tempered &
Conditional on Ramanujan &
\Cref{v4-arith:w5-a:rem:non-tempered}. \\
Eisenstein continuous spectrum &
Scope-excluded &
\Cref{v4-arith:w5-a:prop:eisenstein-residual}. \\
Full $\mathrm{PW}(\mathbb{R})$ class &
Open (H2 obstruction) &
\Cref{v4-arith:w5-a:rem:H2-obstruction}. \\
\end{tabular}
\end{center}

\begin{theorem}[wave-5 HP--Li status]
\label{v4-arith:w5-a:thm:wave5-status}
\ClaimStatusProvedHere. HP--Li positivity
(\Cref{v4-arith:rh:thm:VB-implies-F-RH}) closes unconditionally on
the restricted class $\mathrm{PW}^{\mathrm{hol}}_{F}$ of Paley--Wiener
test functions adapted to a holomorphic cuspidal Hecke eigenform $F$
of weight $k\geq 12$, up to the archimedean Tate distribution
$W_{\infty}$. The prime-side Weil contribution is an instance of
Petersson positivity on the tempered cuspidal core via the
Mellin--Rankin--Selberg isometry $\Upsilon_{F}$
(\Cref{v4-arith:w5-a:def:upsilon}, \Cref{v4-arith:w5-a:thm:RS-integral}).
The irreducible obstruction to full HP--Li positivity
(hence full RH) is the archimedean Tate positivity
\Cref{v4-arith:w5-a:conj:archimedean-tate-positivity}, which is
strictly equivalent to RH itself; the arithmetic chiral-algebra
framework does not close this final obstruction.
\end{theorem}

\begin{proof}
Assemble:
\Cref{v4-arith:w5-a:thm:prime-side},
\Cref{v4-arith:w5-a:thm:RS-integral},
\Cref{v4-arith:w5-a:thm:HPLi-restricted},
\Cref{v4-arith:w5-a:prop:archimedean-not-absorbed},
\Cref{v4-arith:w5-a:thm:ATP-equals-weil-positivity}.
The restricted-class closure is unconditional; the archimedean
residual is RH-equivalent by the reformulation
\Cref{v4-arith:w5-a:thm:ATP-equals-weil-positivity}.
\end{proof}

\begin{corollary}[impact on the minimal RH path]
\label{v4-arith:w5-a:cor:minimal-path-impact}
\ClaimStatusProvedHere. The minimal RH path of
\Cref{v4-arith:rh:thm:minimal-path} is updated on the HP--Li input
as follows:
\begin{itemize}
\item \textbf{HP--Li prime-side: closed} on
$\mathrm{PW}^{\mathrm{hol}}_{F}$ via the Mellin--Rankin--Selberg
isometry.
\item \textbf{HP--Li archimedean contribution: open, RH-equivalent}
(\Cref{v4-arith:w5-a:conj:archimedean-tate-positivity}).
\item \textbf{HP--Li extension beyond
$\mathrm{PW}^{\mathrm{hol}}_{F}$: open}, with two routes:
(a) Maass lifts (Ramanujan-conditional), (b) weight $k\to\infty$
limits (requiring control of the Petersson limit measure, a
problem equivalent to Selberg trace).
\end{itemize}
The irreducible wave-5 residual is therefore a single statement:
archimedean Tate positivity $W_{\infty}(f*\widetilde{f})\leq
W_{\mathrm{spect}}(f*\widetilde{f})$, which is itself equivalent to
RH. The arithmetic chiral-algebra framework has relocated
the HP--Li obstruction from the full positive cone to the archimedean
gamma distribution; it has not closed the obstruction.
\end{corollary}

\begin{proof}
Combine \Cref{v4-arith:w5-a:thm:wave5-status} with the minimal-path
theorem \Cref{v4-arith:rh:thm:minimal-path}.
\end{proof}

\begin{remark}[scope of the wave-5 closure]
\label{v4-arith:w5-a:rem:scope-summary}
The wave-5 closure does not prove RH. What it does:
\begin{enumerate}
\item Isolates the prime-side positive-trace content as unconditional
Petersson positivity on the tempered cuspidal core via an explicit,
computable isometry.
\item Names the archimedean residual precisely as the Tate
$\psi$-distribution, whose positivity is strictly equivalent to RH.
\item Converts HP--Li from a single opaque positivity inequality into
a decomposition
$W(f*\widetilde{f})=W_{\mathrm{prime\text{-}spec}}^{+}(f)-W_{\mathrm{arch}}(f)$,
the first term a closed positive form on the Petersson core, the
second a sign-indefinite archimedean distribution.
\item Furnishes an explicit test-function class
$\mathrm{PW}^{\mathrm{hol}}_{F}$ on which the closed part is
unconditional; this class is dense in $\mathrm{PW}(\mathbb{R})$ in
the Schwartz topology and therefore captures the full positivity
question asymptotically.
\end{enumerate}
The closure of HP--Li on the restricted class is the wave-5
deliverable. The closure of HP--Li on the full Paley--Wiener class
is strictly equivalent to RH and is not closed here.
\end{remark}

\section{Summary and Beilinson-Principled Closure}
\label{v4-arith:w5-a:summary}

\begin{enumerate}
\item The Mellin lift
$\Upsilon_{F}\colon\mathrm{PW}^{\mathrm{hol}}_{F}\to\mathcal{V}^{\mathrm{prim}}_{\mathrm{cusp,temp}}$
of \Cref{v4-arith:w5-a:def:upsilon} is an explicit linear map from
Paley--Wiener test data to the tempered cuspidal core of
$\mathcal{V}^{\mathrm{prim}}$, constructed from a holomorphic cuspidal
Hecke eigenform $F$ of weight $k\geq 12$.
\item The Rankin--Selberg integral
\Cref{v4-arith:w5-a:thm:RS-integral} computes
$\langle\Upsilon_{F}(f),\Upsilon_{F}(f)\rangle^{\mathrm{Pet}}$ as a
non-negative integral on the critical line.
\item The prime-side Weil contribution is an instance of Petersson
positivity via the isometry:
$W_{\mathrm{fin}}(f*\widetilde{f})\geq 0$ on
$\mathrm{PW}^{\mathrm{hol}}_{F}$
(\Cref{v4-arith:w5-a:thm:prime-side}).
\item The irreducible obstruction to full HP--Li is the archimedean
Tate distribution $W_{\infty}$; its positivity is strictly
equivalent to Weil positivity and to RH
(\Cref{v4-arith:w5-a:conj:archimedean-tate-positivity},
\Cref{v4-arith:w5-a:thm:ATP-equals-weil-positivity}).
\item The Eisenstein continuous spectrum is scope-excluded by the
density mismatch of \Cref{v4-arith:kms-gns:cor:eisenstein-final}; it
does not contribute to the Li coefficients, which see only the
discrete spectrum.
\item Wave-5 HP--Li status
(\Cref{v4-arith:w5-a:thm:wave5-status}): HP--Li closes
unconditionally on $\mathrm{PW}^{\mathrm{hol}}_{F}$ up to the
archimedean Tate term; full closure is RH-equivalent and not
achieved.
\end{enumerate}

\begin{remark}[Beilinson principle application]
\label{v4-arith:w5-a:rem:beilinson-principle}
The wave-5 deliverable obeys the Beilinson principle: a smaller
closed theorem (restricted-class HP--Li) is preferred to a larger
unclosed claim (full HP--Li, equivalent to RH). The closed portion
is explicit, computable, and unconditional within the restricted
class; the open portion is named precisely as the archimedean Tate
positivity conjecture A-TP. The isometry
$\Upsilon_{F}$ is the mathematical content that the arithmetic
chiral-algebra framework contributes to the HP--Li programme: a
specific, verifiable lift of Paley--Wiener test data into the
tempered cuspidal core of the adelic chiral algebra, whose Petersson
norm is an explicit Rankin--Selberg integral matching the prime-side
Weil distribution exactly, and leaving the archimedean distribution
as the irreducible residual.
\end{remark}

\section{Bibliography of Primary Sources}
\label{v4-arith:w5-a:bibliography}

The following primary sources underwrite the identities of this
chapter and are disjoint from the derivation sources used in
Chapters~\ref{v4-arith:closure} (P1--P3 resolution),
\ref{v4-arith:kp-compat} (Koszul--Pet compatibility), and
\ref{v4-arith:kms-gns-full} (KMS${}_{1}$/GNS full sector).

\begin{description}
\item[Weil 1952.] Weil, A. \emph{Sur les formules explicites de la
th\'eorie des nombres premiers}. \emph{Meddelanden fr\r{a}n Lunds
Univ.~Mat.~Sem.} (Festschrift f\"or Marcel Riesz), 1952; reprinted
\emph{Oeuvres Scientifiques} Vol.~II, Springer 1979. Explicit formula,
$W$-distribution, Weil positivity equivalence to RH.
\item[Li 1997.] Li, X.-J. \emph{The positivity of a sequence of
numbers and the Riemann hypothesis}. \emph{J.~Number Theory} 65
(1997), 325--333, arXiv:math/9712243. Li coefficients $\lambda_{n}$,
equivalence to RH, explicit Paley--Wiener realisation.
\item[Rankin 1939.] Rankin, R.~A. \emph{Contributions to the theory
of Ramanujan's function $\tau(n)$ and similar arithmetical functions}.
\emph{Proc.~Camb.~Phil.~Soc.}~35 (1939), 357--372. Rankin--Selberg
identity for $L$-series of modular forms.
\item[Selberg 1940.] Selberg, A. \emph{Bemerkungen \"uber eine
Dirichletsche Reihe, die mit der Theorie der Modulformen nahe
verbunden ist}. \emph{Arch.~Math.~Naturvid.}~43 (1940), 47--50.
Rankin--Selberg identity and $L$-function analytic continuation.
\item[Hecke 1936.] Hecke, E. \emph{\"Uber die Bestimmung Dirichletscher
Reihen durch ihre Funktionalgleichung}. \emph{Math.~Ann.} 112
(1936), 664--699. Mellin transform / modular form correspondence
and functional equation.
\item[Deligne 1974.] Deligne, P. \emph{La conjecture de Weil I}.
\emph{Publ.~Math.~IHES} 43 (1974), 273--307. Ramanujan bound for
holomorphic cuspidal Hecke eigenforms of weight $k\geq 2$.
\item[Jacquet 1972.] Jacquet, H. \emph{Automorphic forms on $GL_{2}$
II}. Lecture Notes in Mathematics 278, Springer 1972. Local
archimedean factors for $GL_{2}$; ramified principal series.
\item[Jacquet--Langlands SLN~114 1970.] Jacquet, H.;
Langlands, R.~P. \emph{Automorphic forms on $GL(2)$}. Lecture Notes
in Mathematics 114, Springer 1970. Adelic Whittaker model, Hecke
theory.
\item[Iwaniec--Kowalski 2004.] Iwaniec, H.; Kowalski, E.
\emph{Analytic Number Theory}. AMS Colloquium Publications 53, 2004.
Modern account of Weil's explicit formula, Rankin--Selberg,
Paley--Wiener.
\item[Bombieri--Lagarias 1999.] Bombieri, E.; Lagarias, J.~C.
\emph{Complements to Li's criterion for the Riemann hypothesis}.
\emph{J.~Number Theory} 77 (1999), 274--287. Alternative constructions
of Li-coefficient approximants and equivalences to Weil positivity.
\item[Connes 1999.] Connes, A. \emph{Trace formula in noncommutative
geometry and the zeros of the Riemann zeta function}. \emph{Selecta
Math.~(N.~S.)} 5 (1999), 29--106, arXiv:math/9811068. Adele class
space, absorption spectrum, Weil positivity as trace-formula
positivity.
\item[Patterson 1988.] Patterson, S.~J. \emph{An Introduction to the
Theory of the Riemann Zeta-Function}. Cambridge Studies in Advanced
Mathematics 14, 1988. Complete treatment of the Hadamard product
and explicit formula.
\item[Kim--Sarnak 2003.] Kim, H.~H.; Sarnak, P.~C. Appendix to
Kim, H.~H., \emph{Functoriality for the exterior square of $GL_{4}$
and the symmetric fourth of $GL_{2}$}. \emph{J.~Amer.~Math.~Soc.} 16
(2003), 139--183. Bound $|a_{F}(p)|\leq 2p^{7/64}$ towards the
Ramanujan conjecture for $GL_{2}/\mathbb{Q}$.
\item[Bump 1997.] Bump, D. \emph{Automorphic Forms and
Representations}. Cambridge Studies in Advanced Mathematics 55, 1997.
Comprehensive account of Rankin--Selberg, Whittaker models,
Petersson theory.
\item[Zagier 1981.] Zagier, D. \emph{The Rankin--Selberg method for
automorphic functions which are not of rapid decay}.
\emph{Bonner Math.~Schriften} 121 (1981). Rankin--Selberg unfolding
technique.
\end{description}

\noindent\emph{Disjointness audit.} The primary-source catalog
driving this chapter (Weil 1952 for the explicit formula; Li 1997 for
Li coefficients; Rankin 1939 and Selberg 1940 for the Rankin--Selberg
identity; Hecke 1936 for Mellin--functional-equation duality;
Deligne 1974 for temperedness; Jacquet 1972 and Jacquet--Langlands
SLN~114 for adelic $GL_{2}$ representation theory) is disjoint from
the catalogs underwriting P1 (Beilinson--Drinfeld 2004, Positselski
2009), P2 (Tate 1950), P3 / Koszul--Pet (Bost--Connes 1995,
Takesaki 1979, Jacquet--Piatetski-Shapiro--Shalika 1983,
Godement--Jacquet SLN~260, Moeglin--Waldspurger 1995, Arthur 1980,
Shahidi 2010), the Tate restricted tensor chapter (Lurie HA,
Fresse~2017, Ayala--Francis), and the HP$^{\mathrm{Ar}}$ / VB$^{*}$
converse catalogs (Connes 1999 is a cross-reference cited for
context in both chapters; it does not appear on both sides of any
single identification). The HZ-IV disjointness check for this chapter
passes.
